\chapter{Problem statement and requirements}\label{ch:problem}

\section{The problem}
A developer tunes a query on a laptop or staging server. Production has different CPUs, memory
bandwidth, core counts, storage and often a different instruction set. The question ``will this meet
the latency budget there?'' is normally answered by deploying and measuring, which is slow and
sometimes dangerous. CETP tries to answer it before deployment, using only information obtainable
without privileged access to production.

\section{Why it is hard}
\begin{enumerate}
\item \textbf{A query is not one thing.} Its runtime mixes sequential scanning (bandwidth-bound), hash
  building and sorting (cache- and compute-bound), nested loops (latency-bound) and parallel workers.
  Machines differ on each of these by different factors (serial-axis spread up to $1.39\times$ on the Docker
  grid versus $3.57\times$ for parallelism, \evFILE\ \src{results/operator\_model\_metrics.json}).
\item \textbf{Hardware counters are unavailable.} The AWS Nitro hypervisor blocks the PMU, so the
  well-known ``counter-based prediction'' route is closed for the target cloud \evCLAIM.
\item \textbf{Noise is large relative to the signal.} The first repetition of a cell is the slowest in
  61 of 105 cells, worst case $8.33\times$; after discarding it, 11 of 105 cells still exceed 5\% CV
  (\evFILE\ \src{results/canonical\_metrics.json}). Error numbers below about 5\% are meaningless.
\item \textbf{The planner can change its mind.} Different versions or statistics can choose a different
  plan; 12 of 22 TPC-H plans change across PostgreSQL 14/15/17/18 (research pack, \evPACK).
\item \textbf{Only a few machines exist.} Five EC2 types and eight containers are tiny samples of the
  population of machines, so any learned model is fitted on very few distinct hardware points.
\end{enumerate}

\section{Scope: what is in and out}
\begin{center}\small
\begin{tabularx}{\textwidth}{@{}YY@{}}
\toprule \thead{In scope} & \thead{Out of scope (stated or observed)}\\ \midrule
Predict runtime \emph{scaling} for a query whose dev plan is known & Predict plan changes; the model assumes the same plan\\
Analyse EXPLAIN output; run a probe that creates and drops its own \texttt{UNLOGGED} tables & Touching user data. Note the probe is not read-only: it needs \texttt{CREATE} privilege and about 8.4 million scratch rows across three tables (2.0M + 0.4M + 6.0M)\\
Confidence and refusal when extrapolating (low confidence, cross-ISA) & Absolute-time models of a machine unseen on both sides\\
Concurrency scaling via a simple power law $N^{e}$ & Lock contention, buffer-pool eviction, I/O-bound workloads (classifier recall for I/O is 0)\\
TPC-H scale factor 1 as the benchmark & Real customer workloads (none were measured)\\
\bottomrule
\end{tabularx}
\end{center}

\section{Requirements}
These are reconstructed from the repository (README, CLI, tests, research pack). The original capstone
requirement document was not provided; where a requirement is inferred from code rather than a stated
document, it is marked. IDs are used by the traceability matrix in \cref{app:trace}.

\subsection{Functional requirements}
\begin{center}\small
\begin{longtable}{@{}L{1.25cm}L{8.6cm}L{4.9cm}@{}}
\toprule \thead{ID} & \thead{Requirement} & \thead{Source of requirement}\\ \midrule \endhead
FR-1 & Capture a hardware signature of a machine (bandwidth, compute; operator rates) & README; \src{scripts/collection/operator\_probe.py}\\
FR-2 & Decompose an EXPLAIN ANALYZE plan into per-operator self-times & \src{cetp/plan\_nodes.py}; tests\\
FR-3 & Predict a scaling factor and runtime from dev plan + two signatures & \src{cetp/cross\_env.py}\\
FR-4 & Refuse or warn when extrapolating (cross-ISA uncalibrated, low confidence); \textbf{partially met}: the CLI path bypasses the refusal (\cref{sec:cli-bypass}) & \src{cetp/cross\_env.py}; tests\\
FR-5 & Calibrate from $k\ge3$ production observations; persist atomically & \src{cetp/calibration\_store.py}\\
FR-6 & Apply a concurrency factor $N^{e}$ with elasticity $e$ & \src{cetp/operator\_model.py}\\
FR-7 & Gate a query against an SLA with PASS/WARN/BLOCK & \src{gate/cetp\_gate.py}\\
FR-8 & Detect regressions between two plans on the same hardware & \src{cetp/regress.py}\\
FR-9 & Ingest queries from \texttt{auto\_explain}/\texttt{pg\_stat\_statements} & \src{cetp/ingest.py}\\
FR-10 & Redact sensitive literals/identifiers before sharing plans & \src{cetp/redact.py}\\
FR-11 & Local telemetry without query text or identifiers & \src{cetp/telemetry.py}\\
FR-12 & Command-line interface for all of the above & \src{cetp/cli.py}\\
\bottomrule
\end{longtable}
\end{center}

\subsection{Non-functional requirements}
\begin{center}\small
\begin{tabularx}{\textwidth}{@{}L{1.25cm}Y L{4.3cm}@{}}
\toprule \thead{ID} & \thead{Requirement} & \thead{Where checked}\\ \midrule
NFR-1 & Reproducibility: published numbers regenerate byte-for-byte & \texttt{make verify-determinism} (\evPACK)\\
NFR-2 & No runtime dependencies; Python $\ge$ 3.11 & \src{pyproject.toml}\\
NFR-3 & Leakage-free evaluation (no test-fold information in fitting) & Leakage audit in \src{canonical\_evaluation.py}\\
NFR-4 & Honest uncertainty: report CIs, paired differences, and negative results & \src{results/CANONICAL.md}\\
NFR-5 & Privacy: telemetry and shared plans contain no query text/identifiers & \src{cetp/redact.py}, \src{cetp/telemetry.py}\\
NFR-6 & Crash-safe persistence of calibration & atomic write in \src{cetp/calibration\_store.py}\\
NFR-7 & Probe must not leave artefacts on the target database & \texttt{finally}-block drops in \src{operator\_probe.py}\\
\bottomrule
\end{tabularx}
\end{center}

\subsection{Success criteria}
The repository states its own bar implicitly: a predictor is ``useful'' only if it beats the naive
baseline with a confidence interval that excludes zero. By that bar, the status per scenario is:
\begin{center}\small
\begin{tabularx}{\textwidth}{@{}L{4.2cm}L{2.2cm}L{2.2cm}Y@{}}
\toprule \thead{Scenario} & \thead{Best model} & \thead{Naive} & \thead{Verdict}\\ \midrule
EC2 LOMO (840 pairs) & 21.61 & 31.96 & Met vs naive; not vs compute-only\\
Docker grid (1{,}232 pairs) & 22.41 & 72.52 & Met, large margin\\
Laptop $\to$ prod, same ISA (7 cells) & 29.31 & 37.39 & Met\\
Laptop $\to$ prod, all 11 cells & 43.72 & 31.70 & \textbf{Not met}\\
Laptop $\to$ prod, x86 runners (4) & 68.92 & 21.74 & \textbf{Not met} (3$\times$ worse)\\
Laptop $\to$ EC2 (95 pairs) & 160.68 & 129.39 & \textbf{Not met}\\
\bottomrule
\end{tabularx}
\end{center}
All values are MAPE (\%), \evFILE; sources in \cref{app:evidence}.

\section{The central trade-off}
A predictor can be safe by being right more often or by being honest about when it is not. CETP does
the second at best: confidence-aware abstention helps modestly: dropping the least confident 20\% of predictions leaves a
selective MAPE of 28.13\%, and the area under the risk--coverage curve is 25.44 for hardware-distance
ordering against 31.15 for random ordering (\evFILE\ \src{results/confidence\_abstention\_summary\_corrected.json}), but
this does not turn an inaccurate predictor into an accurate one.
